\section{The Ising pre-decoders and their evaluation}\label{sec:background}

\textbf{Input window.} The pre-decoder sees a window of $T$ syndrome rounds of a rotated surface code of distance $d$
as a tensor of shape $(4, T, d, d)$ \cite{chamberland2026predecoder}. Two channels carry the syndrome bits of the X
and Z stabilizers, placed on the $d\times d$ grid of data qubits, and two channels carry fixed masks that mark which
grid positions hold a stabilizer of each type. The syndrome channels take the values 0 and 1 and the mask channels
the values 0, $\tfrac12$ and 1.

\textbf{Network.} Both released surface code networks are stacks of three dimensional convolutions with $3\times
3\times 3$ kernels, zero padding that keeps the window size, and the tanh approximation of GELU
\cite{hendrycks2016gelu} between layers, without normalization layers. The Fast network (model 1) has four layers,
$4\to128\to128\to128\to4$ channels and 912{,}772 parameters; the Accurate network (model 4) has six layers,
$4\to128\to128\to128\to128\to128\to4$ channels and about 1.8 million parameters \cite{nvidia_fast_card,
nvidia_accurate_card}. Their receptive fields are 9 and 13 sites, and they were trained on windows with $d=T=9$ and
$d=T=13$. Because the networks are fully convolutional, the same weights apply to any larger window.

\textbf{Output and use.} The four output channels are logits of a correction on each data qubit (two channels) and
of a flip of each syndrome bit (two channels). A logit at or above a threshold (zero in the released configuration)
selects the correction. The corrections are applied to the syndrome, and PyMatching decodes the residual syndrome
with the detector error model of the circuit; the logical outcome combines the pre-decoder corrections and the
matching \cite{chamberland2026predecoder,higgott2025sparse}.

\textbf{Evaluation.} NVIDIA's evaluation samples memory experiments with Stim \cite{gidney2021stim} under a circuit
level noise model with 25 parameters (state preparation, measurement, idling during CNOT layers and during
preparation and reset, and the fifteen two qubit Pauli errors of a CNOT), which their public configuration sets for a
physical error rate $p=0.003$ \cite{nvidia_ising_code}. For each shot it reports whether PyMatching alone and the
pre-decoder followed by PyMatching recover the logical observable. We use this code without modification for every
logical error rate in this paper, and add only the scaling of the 25 parameters to other values of $p$ with the
relations given in the public configuration (preparation and measurement $2p/3$, idling during CNOT layers $p/3$ per
Pauli, idling during preparation and reset $2p/3-4p^2/9$ per Pauli, and $p/15$ per two qubit Pauli). We report the
ratio of the number of logical errors of PyMatching alone to that of the pre-decoder followed by PyMatching on the
same shots; a ratio above one is an improvement.
